\section{Introduction}

An agent that quickly adapts to a large variety of tasks in an environment needs to capture a useful representation which encodes how its behavior affects the environment and its own state within it.
To accomplish this, various unsupervised reinforcement learning techniques such as goal-conditioned reinforcement learning, unsupervised skill learning, forward-backward representation learning, and proto successor measure approximations have been proposed.
Recently, \citet{agarwal-rlj25a} have shown that these objectives can be unified as approximations to the successor measure, a generalization of the future state-action occupancy of an agent.

Explicit successor measure approximations such as proto successor measures~\citep{agarwal-arxiv25a} and forward-backward representations~\citep{touati-neurips21a,touati2023does,trinzoni-iclr25a,bagatella-iclr26a} need to be defined with regard to a set of basis policies or tasks.
These prior approaches have explicitly sought to obtain representations capable of representing the occupancy under all possible policies.
However, capturing the behavior of all possible policies from a static dataset is in many cases intractable as the amount of data necessary to accomplish this is simply too large.
Some practical approaches have therefore added ad-hoc regularization to the existing behavior foundation model framework, such as behavior cloning losses. %%AZ.8.29: citation? 

 %%AZ.8.29: Need a link here explaining what's wrong with the above approaches, why should someone prefer the below.
We establish a principled way of extracting a diverse set of well-covered basis policies from a static dataset and obtain a successor measure approximation for this dataset.
For this, we use an expressive flow-based policy to capture the dataset distribution. 
Following \citet{wagenmaker}, we treat the input noise of the flow model as a latent action, and define deterministic basis policies in this latent space.
We can then treat this noise space as an unconstrained action space and any policy defined over it will stay close to the set of policies well-covered by the data.

 %%AZ.8.29: the contribution of this paragraph is not coming through. Is the point that we are trying to augment the latent action dataset? What is the benefit of doing this?
However, while flow policies allow us to define unconstrained basis policies, the data itself contains environment actions, not latent actions.
Multiple latent actions can lead to the same environment action, and identifying these allows us to augment the dataset with all valid latent transitions.
We therefore present an algorithm to identify a distribution over latent noise actions which are decoded by the flow policy to similar environment action.
This technique allows us to augment the training dataset for the proto successor measures, but we believe it is also of general interest to the community for many other applications.
To construct the noise space sets, we leverage the local invertibility of the flow model.
We then compute a proposal distribution in the latent space and treat the set identification as a posterior inference problem.

The resulting algorithm for successor measure approximation using inverted flow models and latent actions is called Proto Successor Measure Flows, a principled method for computing successor approximations over limited coverage datasets.
We implement and validate our technique on the OGBench tasks, a goal-conditioned environment and dataset that naturally allows for testing multi-task goal-reaching capabilities of behavior foundation models.
We show ... \todo[inline]{add results here}.